\documentclass[10pt,a4paper]{amsart}
\usepackage[margin=19mm]{geometry}
\usepackage{amsmath,amssymb,mathtools}
\usepackage[scr=rsfs,cal=euler]{mathalfa}
\usepackage{xfrac}
\usepackage[unicode,hidelinks,bookmarks=false]{hyperref}
\newcommand{\ie}{i.e.\ }
\newcommand{\cf}{cf.\ }
\newcommand{\resp}{respectively\ }
\newcommand{\Subsection}{\subsection}
\providecommand{\nobreakdash}{\nobreak\hbox{-}}
\setlength{\emergencystretch}{2em}
\multlinegap0pt
\usepackage{amssymb}
\usepackage{bm}
\usepackage{braket}
\usepackage{xcolor}
\usepackage{tikz}
\newtheorem{lemma}{Lemma}
\newcommand{\End}{\operatorname{End}}
\newcommand\Tr{\operatorname{Tr}}
\newcommand{\cone}{(1)}
\newcommand{\ctwo}{(2)}
\newcommand{\id}{\operatorname{id}}
\title{\fontsize{15pt}{15pt}Weak Hopf non-invertible symmetry-protected topological spin liquid and lattice realization of (1+1)D symmetry topological field theory}
\author{Zhian Jia}
\date{Source: arXiv:2412.15336v2 (CC BY 4.0)}
\begin{document}
\maketitle

\noindent\emph{Source and licence.} \fontsize{15pt}{15pt}Weak Hopf non-invertible symmetry-protected topological spin liquid and lattice realization of (1+1)D symmetry topological field theory. Version arXiv:2412.15336v2 (CC BY 4.0), distributed by arXiv under the Creative Commons Attribution 4.0 International licence, http://creativecommons.org/licenses/by/4.0/, as recorded in the arXiv licence metadata for this exact version and shown on the abstract page https://arxiv.org/abs/2412.15336v2. Full licence notice: \texttt{licenses/arxiv-2412.15336-CC-BY-4.0.txt}.\par\smallskip

\noindent\textit{Editorial scope.} Counted unit: the article's lemma prop:orthgonality (source0.tex lines 1310--1321). The article's complete original proof of it is delivered with it (source0.tex lines 1324--1338, one \texttt{proof} environment of 15 lines). No mathematics is written, repaired or re-derived: the statement and the proof are the article's own lines, in the article's own notation, and no retained line is substituted or re-flowed.  Retained uncounted as the source-local context the proof needs: lines 1287--1289; lines 1301--1303.  Nothing was omitted from any retained span, so the package carries no omission record.  No retained line contains a citation, so the wrapper carries no bibliography and no bibliography entry.  No retained line contains a figure environment, an image-inclusion command or drawing code, so no image is delivered and the package declares no asset dependency.  Editorial formatting only, in addition: an AMS \texttt{amsart} preamble that restates the article's own declarations and macros used by the retained lines, namely the environments \texttt{lemma} on separate counters, together with the standard packages those lines need.  The retained environments print in this document as Lemma 1 in order of appearance, whereas the article prints the same items with the numbering of its own full document.  The complete native source of the article as distributed in the arXiv e-print for this exact version is bundled unchanged as \texttt{sources/170439/source0.tex}.
\begin{lemma}[Schur's lemma for weak Hopf algebra]\label{prop:schur} 
Let $\Gamma: H \to \End(V)$ and $\Phi: H \to \End(W)$ be two irreducible representations of a weak Hopf algebra $H$, and let $f: V \to W$ be a linear map such that $\Phi(x) \circ f = f \circ \Gamma(x)$ for all $x \in H$. If $\Gamma$ and $\Phi$ are not isomorphic, then $f = 0$. When $\Gamma= \Phi$, $f$ is a scalar multiple of the identity map on $V$, i.e., $f \propto \id_V$. 
\end{lemma}

\begin{equation}\label{eq:Fave}
    F=\sum_{(\lambda)} \Phi(S(\lambda^{(1)})) \circ f \circ \Gamma(\lambda^{\ctwo}).
\end{equation}

\begin{lemma}\label{prop:orthgonality}
 (1)  For any $f:V\to W$, and define $F$ as in Eq.~\eqref{eq:Fave}, we have: (i)  If $\Gamma$ and $\Phi$ are not isomorphic, $F=0$; (ii) when $\Gamma=\Phi$, we have
    \begin{equation}\label{eq:Ftrace}
        F=\frac{1}{d_{\Gamma}}\Tr (f) \id_V.
    \end{equation}
    
(2) From the above result, we obtain the orthogonality relation for irreps of $H$:
\begin{equation}
    \sum_{(\lambda)} \Phi_{ij}(S(\lambda^{\cone})) \Gamma_{kl}(\lambda^{\ctwo})=\frac{\delta_{\Phi,\Gamma} \delta_{il}\delta_{jk} }{d_{\Gamma}}=  \sum_{(\lambda)} \Phi_{ij}(S(\lambda^{\ctwo})) \Gamma_{kl}(\lambda^{\cone}).
\end{equation}
where we have used the cocommutativity of $\lambda$.
\end{lemma}

\begin{proof}
   (1) Note that (i) is a direct result of Proposition~\ref{prop:schur}.  For (ii), we have
   \begin{equation}
        \Tr F= \sum_{(\lambda)} \Tr [\Gamma(S(\lambda^{(1)})) \circ f \circ \Gamma(\lambda^{\ctwo})]=\Tr [\Gamma(\varepsilon_L(\lambda)) f]=\Tr f,
    \end{equation}
which implies Eq.~\eqref{eq:Ftrace}, notice that we have used $\varepsilon_L(\lambda)=1_H$.

(2) From (1), the matrix elements $F_{il}$ is of the form
\begin{equation}
    \begin{aligned}
         \sum_{(\lambda)} \sum_{jk} \Phi_{ij}(S(\lambda^{\cone}))f_{jk} \Gamma_{kl}(\lambda^{\ctwo})=  \frac{\delta_{\Phi,\Gamma}}{d_{\Gamma}} \delta_{il} \sum_{jk} f_{jk}\delta_{jk}.
    \end{aligned}
\end{equation}
This implies the required result by taking $f_{jk}=\delta_{jj'}\delta_{kk'}$.
\end{proof}

\end{document}
